% Desarrollo reunido K31-07--K31-15. Inserción anterior a registro_k y alpha.
% Mantiene explícita la interfaz del repertorio terminal de ocho bloques.
% Propietarios y localizadores: metadata/K_ALPHA_PROPIETARIOS.md.
\section{Determinación regional del descriptor dodecafásico y selección del registro de acoplamiento}
\label{act21:chap:despliegue-selector-k}

La determinación del registro de acoplamiento articula dos construcciones
que conviene desplegar en su orden efectivo. La primera produce los
antecedentes del selector a partir de las representaciones regionales del TPK:
bandas de seis trits, firmas de columna, calendario de conversión,
descriptor orientado y panel funcional. La segunda opera sobre esos
antecedentes y sobre el repertorio terminal de ocho bloques, enumera sus
ordenaciones admisibles y selecciona un único registro mediante dos
condiciones simultáneas: incidencia excepcional y heterotipia de lectura.
El valor del registro aparece al término de esta composición.

El presente desarrollo presupone las representaciones regionales obtenidas
en los capítulos precedentes desde APP--TRIT--TPK. Su propósito es hacer
explícita cada operación intermedia que conduce desde tales representaciones
al selector finito. Después se incorporan la realización del mismo
registro por incidencias firmadas, la inversión de Hadamard y la
representación de la constante de estructura fina. Las tres construcciones
conservan sus respectivos dominios: selección del registro, recuperación
desde sus canales y evaluación aritmética de sus representaciones.

\subsection{Bandas regionales, codificación y firmas enteras}

\subsubsection{Dominios y codificación de las bandas regionales}

Se fija el orden de canales
\[
 \chi\in\{\pi,e,\varphi\},
 \qquad B_\chi(t)\in\mathbb F_3^6,
 \qquad t\in\mathbb N.
\]
Los símbolos de canal identifican las regiones ya generadas. Cada
$B_\chi(t)$ conserva el lugar que ocupa en su sucesión de representaciones;
el índice $t$ cuenta bloques ternarios. La evaluación natural de una
banda $b=(b_0,\ldots,b_5)$ es
\begin{equation}
 \operatorname{enc}_6(b)=
 243b_0+81b_1+27b_2+9b_3+3b_4+b_5,
 \qquad 0\leq\operatorname{enc}_6(b)<729.
 \label{act21:eq:kdes-enc6}
\end{equation}
Su inversa, con anchura fijada, es
\[
 \operatorname{dig}_j(n)=
 \left\lfloor\frac{n}{3^{5-j}}\right\rfloor\bmod3,
 \qquad 0\leq j<6.
\]
Así se retienen los ceros iniciales: las palabras $002001$ y $2001$
podrían compartir una escritura entera abreviada, pero el tipo de la
primera contiene seis posiciones. Todas las concatenaciones que siguen
utilizan esa anchura.

En la ventana de cien bloques, la representación de seiscientos trits del
canal $\chi$ determina un entero $P_\chi^{(600)}$. La extracción se efectúa
mediante
\begin{equation}
 b_\chi(t)=
 \left\lfloor\frac{P_\chi^{(600)}}{729^{99-t}}\right\rfloor\bmod729,
 \qquad 0\leq t<100,
 \qquad B_\chi(t)=\operatorname{enc}_6^{-1}(b_\chi(t)).
 \label{act21:eq:kdes-extraccion}
\end{equation}
El entero de seiscientos trits es una representación del productor regional;
la tabla de cien filas es la evaluación de
\eqref{act21:eq:kdes-extraccion}. Esta separación permite sustituir una tabla
histórica por su productor sin modificar el selector que la consume.

\subsubsection{Residuos, cocientes, ocupación y orientación}

Para cada columna $j$ y tiempo $t$, se introducen
\begin{align}
 T_j(t)&=B_\pi(t)_j+B_e(t)_j+B_\varphi(t)_j,\\
 q_j(t)&=T_j(t)\bmod3,
 &c_j(t)&=\left\lfloor T_j(t)/3\right\rfloor,\\
 a_j(t)&=\bigl(B_\pi(t)_j+B_e(t)_j-B_\varphi(t)_j\bigr)\bmod3,\\
 w_j(t)&=\mathbf1_{B_\pi(t)_j\ne0}
       +\mathbf1_{B_e(t)_j\ne0}
       +\mathbf1_{B_\varphi(t)_j\ne0},
 &\bar w_j(t)&=w_j(t)\bmod3.
 \label{act21:eq:kdes-firmas}
\end{align}
La resta de la tercera expresión se realiza en los enteros y se reduce
después. En una implementación sobre naturales, la expresión equivalente
es $(B_\pi+B_e+3-B_\varphi)\bmod3$; la resta truncada alteraría el lector.

\begin{proposition}[Reconstrucción de la suma de columna]
Para todas las bandas admisibles,
\[
 T_j(t)=q_j(t)+3c_j(t),\qquad
 0\leq q_j(t)<3,\quad 0\leq c_j(t)\leq2,
 \quad0\leq w_j(t)\leq3.
\]
\end{proposition}
\begin{proof}
Cada columna suma tres representantes de $\{0,1,2\}$, de modo que
$0\leq T_j(t)\leq6$. La división euclídea por tres proporciona el residuo
y el cociente indicados. La ocupación cuenta tres condiciones binarias;
su cota se deduce directamente de esa definición.
\end{proof}

Esta identidad conserva el cociente de la suma local. La reconstrucción
de las tres bandas individuales requiere las bandas mismas o un lector
adicional que las distinga: una suma de columna tiene menos coordenadas
que el triple que la produce. Por ello, la fila temporal utilizada más
adelante conserva conjuntamente las tres bandas y los cuatro campos
$q,a,c,\bar w$.

\subsection{Reloj de conversión y primer avance nulo}

\subsubsection{Capacidad ternaria y profundidad decimal}

Una representación de $t+1$ bandas contiene $6(t+1)$ trits. El calendario
de conversión asociado se define por el entero
\begin{equation}
 d_t=\max\{k\in\mathbb N:1000^k\leq3^{6(t+1)}\}
     =\max\{k\in\mathbb N:1000^k\leq729^{t+1}\}.
 \label{act21:eq:kdes-reloj}
\end{equation}
La definición compara capacidades enteras. El contador $t$ y la
profundidad $d_t$ registran magnitudes distintas: una transición añade
una banda de seis trits; el reloj cuenta potencias completas de la base
decimal mil contenidas en la capacidad ternaria acumulada.

\begin{lemma}[Caracterización entera del reloj]
Para $t,k\in\mathbb N$ se tiene
\[
 d_t=k\quad\Longleftrightarrow\quad
 1000^k\leq729^{t+1}<1000^{k+1}.
\]
Además, $d_t\leq d_{t+1}$.
\end{lemma}
\begin{proof}
Las potencias de mil forman una sucesión estrictamente creciente y no
acotada. La potencia $729^{t+1}$ queda entre dos términos consecutivos;
el exponente inferior es exactamente el máximo de
\eqref{act21:eq:kdes-reloj}. La desigualdad
$729^{t+1}<729^{t+2}$ implica la monotonía del máximo.
\end{proof}

\begin{proposition}[Primer avance nulo]
El menor tiempo que satisface $d_{t+1}=d_t$ es
\[
 t_*=20.
\]
En particular,
\[
 d_0=0,\quad d_{19}=19,\quad d_{20}=20,
 \quad d_{21}=20,\quad d_{22}=21.
\]
\end{proposition}
\begin{proof}
Las comparaciones enteras exactas
\[
 1000^{20}\leq729^{21},\qquad
 729^{22}<1000^{21},\qquad
 1000^{21}\leq729^{23}<1000^{22}
\]
dan los valores en los tiempos veinte, veintiuno y veintidós. Para
$0\leq t\leq20$, la expresión
$729^{t+1}/1000^t=729(729/1000)^t$ es decreciente y su valor en veinte
es mayor o igual que uno. Por tanto,
$1000^t\leq729^{t+1}<1000^{t+1}$ y $d_t=t$. Ningún tiempo menor que
veinte puede satisfacer $d_{t+1}=d_t$; en veinte sí se cumple.
Todas las comparaciones utilizadas son de potencias enteras.
\end{proof}

El número veinte es, por tanto, una salida del criterio mínimo. La regla
que selecciona el instante es ``primer avance nulo del reloj''; el
algoritmo y su prueba pueden emplear esta regla sin recibir veinte como
parámetro externo.

\subsubsection{Retorno nonádico y antecedente temporal}

La longitud del retorno de fase determina
\begin{equation}
 r=t_*-9=11.
 \label{act21:eq:kdes-retorno}
\end{equation}
Las consultas del descriptor se efectúan en $t_*+1$, $r+1$ y $r$.
Esta distribución conserva la relación entre el instante de pausa, la
traslación nonádica y la lectura de fase inmediatamente asociada. Los
tiempos once, doce y veintiuno designan posiciones del calendario
regional; su diferencia temporal permanece registrada aun cuando dos
tiempos produzcan la misma profundidad decimal.

\subsection{Construcción orientada del descriptor de veinticuatro trits}

\subsubsection{Defectos opuestos de semibanda}

Se toman los defectos orientados de la coordenatización regional
\[
 \delta_+=(0,0,0,1,1,1),\qquad
 \delta_-=(0,0,0,2,2,2)\in\mathbb F_3^6.
\]
Su oposición se verifica componente a componente:
\[
 \delta_++\delta_-=0\quad\text{en }\mathbb F_3^6.
\]
El representante $2$ realiza $-1$ en la coordenatización ternaria. La concatenación
de las colas de tres posiciones, en el orden negativo--positivo, produce
la palabra $222111$. La suma de defectos y la concatenación de sus
colas son operaciones distintas: la primera tiene codominio
$\mathbb F_3^6$, mientras que la segunda conserva el orden de dos
segmentos de longitud tres.

\subsubsection{Carga dual y selección de fase}

En la coordenatización regional, la transformación de seis coordenadas utiliza
\[
 A_{\rm reg}=
 \begin{pmatrix}
 0&1&1&1&1&1\\
 1&0&1&1&2&2\\
 1&1&0&2&1&2\\
 1&1&2&0&2&1\\
 1&2&1&2&0&1\\
 1&2&2&1&1&0
 \end{pmatrix}.
\]
Con vectores fila, la carga dual de $b\in\mathbb F_3^6$ es la suma
de las coordenadas de $bA_{\rm reg}$. Las sumas enteras de las filas de
$A_{\rm reg}$ son $(5,7,7,7,7,7)$; de ahí resulta la expresión local
\begin{equation}
 \eta^\vee(b)=2b_0+b_1+b_2+b_3+b_4+b_5\pmod3.
 \label{act21:eq:kdes-cargadual}
\end{equation}
La fase de lectura en el instante $r$ es
$\theta_r=(r-1)\bmod3$. Para cada canal se define
\[
 \varepsilon_\chi(r)=
 \begin{cases}
 1,&\eta^\vee(B_\chi(r))=\theta_r,\\
 2,&\eta^\vee(B_\chi(r))\ne\theta_r.
 \end{cases}
\]
La combinación de fase y carga es
\begin{equation}
 D(r)=\varepsilon_\pi(r)B_\pi(r)
      +\varepsilon_e(r)B_e(r)
      +\varepsilon_\varphi(r)B_\varphi(r)
      \quad\text{en }\mathbb F_3^6.
 \label{act21:eq:kdes-fasecarga}
\end{equation}
El coeficiente se determina por una igualdad de cargas, antes de evaluar
la palabra que resultará de la combinación.

\subsubsection{Composición del descriptor}

Se denota por $\Vert$ la concatenación de palabras de anchura fijada.
La construcción completa es
\begin{equation}
 \begin{split}
 \Omega_{24}={}&
 (B_e(t_*+1)+\delta_-)
 \Vert (\delta_-|_{3,4,5}\Vert\delta_+|_{3,4,5})\\
 &\Vert(B_\pi(r+1)+B_e(r+1))\Vert D(r).
 \end{split}
 \label{act21:eq:kdes-omega}
\end{equation}
Cada sumando de longitud seis se evalúa en $\mathbb F_3^6$; la
concatenación final tiene longitud $6+6+6+6=24$. Esta construcción
separa el descriptor $\Omega_{24}$ de las etapas regionales
$w_{24}^{\chi}$: ambos objetos comparten longitud, pero sus productores
y sus lugares en la composición son diferentes.

\begin{proposition}[Evaluación regional del descriptor]
Las representaciones regionales y el criterio del primer avance nulo
determinan
\[
 \Omega_{24}=020022\,222111\,211201\,021101,
 \qquad [\Omega_{24}]_3=66241910521.
\]
\end{proposition}
\begin{proof}
Las bandas necesarias, conservando el orden $(\pi,e,\varphi)$, son
\[
 \begin{array}{c|ccc}
 t&B_\pi(t)&B_e(t)&B_\varphi(t)\\\hline
 11&022212&110001&100021\\
 12&012012&202222&000120\\
 21&221022&020100&211021
 \end{array}
\]
y proceden de la extracción \eqref{act21:eq:kdes-extraccion}. En el tiempo
veintiuno,
$020100+000222=020022$ en $\mathbb F_3^6$. Las colas de los defectos
aportan $222111$. En el tiempo doce,
$012012+202222=211201$.

Para el tiempo once, las cargas duales son, respectivamente, $0,1,2$;
la fase vale $(11-1)\bmod3=1$. Los coeficientes resultan $(2,1,2)$ y
\[
 2(022212)+(110001)+2(100021)=021101
 \quad\text{en }\mathbb F_3^6.
\]
La concatenación de estos cuatro resultados da la palabra anunciada.
Su evaluación entera se obtiene por Horner en base tres, o por tres
concatenaciones sucesivas en base $729$.
\end{proof}

La tabla anterior desempeña una función comprobatoria: cada una de sus
entradas se obtiene del productor regional. La regla de construcción
\eqref{act21:eq:kdes-omega} es la que permite reproducir el descriptor sin
introducir como dato la cadena final de veinticuatro trits.

\subsection{Panel funcional y repertorio terminal}

\subsubsection{Nueve posiciones regionales}

El panel funcional conserva las tres primeras bandas de cada uno de los
tres canales, en el orden regional declarado:
\begin{equation}
 \mathcal P_9=
 \bigl(b_\pi(0),b_\pi(1),b_\pi(2),
       b_e(0),b_e(1),b_e(2),
       b_\varphi(0),b_\varphi(1),b_\varphi(2)\bigr).
 \label{act21:eq:kdes-panel}
\end{equation}
Su evaluación es
\[
 \mathcal P_9=(103,161,103,523,457,301,450,398,438).
\]
Las nueve posiciones conservan procedencia regional, aunque el valor
103 aparece dos veces. El panel tiene nueve incidencias ordenadas y
ocho valores distintos. La enumeración ulterior emplea las incidencias;
la deduplicación de candidatos se efectúa en una etapa posterior.

\subsubsection{Repertorio no ordenado de ocho bloques}

El repertorio terminal que interviene en esta construcción es
\begin{equation}
 \mathcal S_{\rm term}=
 \{002001,020111,200110,121012,
   010122,001001,221110,222110\}\subset\mathbb F_3^6.
 \label{act21:eq:kdes-sterm}
\end{equation}
Sus ocho elementos son distintos. La codificación entera, ordenada por
valor únicamente para fijar una representación informática, es
\[
 \{28,55,98,175,437,498,687,714\}.
\]
La procedencia documental de este repertorio es la segmentación terminal
de la palabra de setenta y dos trits y la recuperación por posiciones
mediante incidencias de los lectores fase--carga y afín. El selector que
se desarrolla a continuación recibe el repertorio sin su ordenación
histórica y somete a prueba sus $8!$ ordenaciones.

La formulación exacta de la implementación reunida mantiene aquí una
interfaz explícita: los productores regionales sustituyen materialmente
el descriptor, el panel y las filas temporales; el repertorio
\eqref{act21:eq:kdes-sterm} conserva su productor documental anterior. Esta
distinción identifica dónde entra cada antecedente. La pertenencia de
ocho palabras a un repertorio temporal de lectores, la selección de esas
ocho palabras y la determinación de su orden terminal son tres
enunciados diferentes. La prueba de selección terminal demuestra el
tercero para el repertorio declarado.

La recuperación documental por posiciones puede expresarse de forma
exacta. Sea $\mathcal I$ el registro de incidencias fase--carga y sea
$\mathcal J$ el registro de la extensión afín. Cada incidencia conserva
la posición $\ell$ de la palabra de setenta y dos trits y la palabra
emitida $u$. Para $5\leq\ell\leq12$ se forma
\[
 \mathcal S_\ell=
 \{u:(\ell,u)\in\mathcal I\cup\mathcal J\}.
\]
El recuperador comprueba que las ocho posiciones aparecen y que cada
$\mathcal S_\ell$ es un singleton. Si $s_\ell$ es su elemento único,
la salida entregada al selector es
\[
 \mathcal S_{\rm term}=\{s_5,s_6,\ldots,s_{12}\}.
\]
La prueba de singleton por posición conserva la coherencia de los
testigos históricos; la exportación posterior conserva las ocho
palabras y suprime únicamente la ordenación histórica. El selector
reconstruye el orden mediante la enumeración completa y las condiciones
terminales, en lugar de recibirlo del registro de procedencia.

% RESULTADO_RECUPERADO: K31-10/P10 del inventario REV05.
% Fuente de las incidencias:
% 16_CIERRE_GLOBAL_HMT_MD_2026-07-22/02_LECTURA_DODECAFASICA/continuacion_fases/
% verificar_obstruccion_y_dato_minimo_w72.py, lineas 45--60 y 141--233.
% RESULTADO_OBSTRUCCION_Y_DATO_MINIMO_W72.json, phase_charge_reader.
% Fuente de las bandas: N69_signatures_t0_t99.csv, tiempos impresos abajo.
% Consumidor: PAQUETE_CONTINUIDAD_K_UNIDAD_20260918_BASE_COMPARTIDA/python/
% selector_orbital_original.py, verify_historical_inputs, lineas 159--193.
% No se transfiere la restriccion del lector comparativo al generador completo.
\subsubsection{Incidencias temporales del repertorio terminal y recuperación por posición}
\label{act21:sec:repertorio-incidencias-completas}

La representación del repertorio terminal utiliza los mismos bloques regionales,
los mismos acarreos y el mismo calendario que producen el descriptor inicial
de veinticuatro trits. Conviene explicitar por separado dos operaciones de
esta composición. La primera recupera ocho palabras a partir de incidencias
temporales registradas. La segunda presenta esas ocho palabras al selector
dodecafásico como un conjunto no ordenado. El selector vuelve a establecer
el orden mediante sus condiciones de incidencia, de cilindro y de cierre.
La distinción permite conservar los testigos de procedencia sin convertir
el orden de una segmentación histórica en una entrada del selector.

Sea $B_t=(b_{t,0},b_{t,1},b_{t,2})\in(\mathbb F_3^6)^3$ la frontera
regional ya generada en el tiempo $t$, con $0\leq t<100$. En la coordenatización
utilizada por el registro histórico, la matriz de incidencia es
\[
 A_{\mathrm{hist}}=
 \begin{pmatrix}
 0&1&1&1&1&1\\
 1&0&1&1&2&2\\
 1&1&0&2&1&2\\
 1&1&2&0&2&1\\
 1&2&1&2&0&1\\
 1&2&2&1&1&0
 \end{pmatrix}.
\]
Las palabras se consideran vectores fila. Sus dos cargas y la fase
cronológica son
\[
 r_{t,i}=\sum_{j=1}^{6}b_{t,i,j},\qquad
 d_{t,i}=\sum_{j=1}^{6}(b_{t,i}A_{\mathrm{hist}})_j,
 \qquad \vartheta_t=t-1\pmod 3.
\]
Todas las operaciones de estas tres expresiones tienen lugar en
$\mathbb F_3$. El subíndice histórico identifica una coordenatización concreta de
la incidencia de Witt; mantiene explícita la transformación de coordenadas
cuando se utiliza otra presentación de esa misma incidencia.

Para $c\in\{r_t,d_t\}$, $\delta\in\mathbb F_3$ y
$\epsilon\in\{1,-1\}$, definimos
\[
 \eta_i(c,\delta,\epsilon)=
 \begin{cases}
  \epsilon,&c_i=\vartheta_t+\delta,\\
  -\epsilon,&c_i\ne\vartheta_t+\delta,
 \end{cases}
 \qquad
 L_{t,c,\delta,\epsilon}
   =\sum_{i=0}^{2}\eta_i(c,\delta,\epsilon)b_{t,i}.
\]
La orientación $-1$ se escribe como $2$ al imprimir los coeficientes
ternarios. Junto a estos doce lectores comparativos por frontera, el
registro dispone del lector afín
\[
 L^{\mathrm{af}}_t
    =\sum_{i=0}^{2}(d_{t,i}+\vartheta_t)b_{t,i}.
\]
La adición de la fase en el lector afín y la comparación de cargas en
el lector anterior son operaciones distintas, ambas sobre datos ya
producidos por la frontera TPK.

\begin{table}[htbp]
\centering\small
\begin{tabular}{rccc cc}
\toprule
$t$&$b_{t,0}$&$b_{t,1}$&$b_{t,2}$&$r_t$&$d_t$\\
\midrule
8 &200121&212212&110110&011&202\\
11&022212&110001&100021&001&012\\
30&210112&110202&211110&100&012\\
34&121022&000011&112110&220&021\\
37&120121&011202&112221&100&201\\
49&110212&200212&100122&110&201\\
58&111121&022121&122201&122&220\\
67&010012&001220&011120&122&122\\
75&220000&020111&100002&120&021\\
81&122201&020202&201122&202&001\\
90&022112&022110&121010&202&200\\
\bottomrule
\end{tabular}
\caption{Fronteras temporales que realizan las incidencias impresas del
repertorio terminal. Cada entrada de seis cifras conserva el orden de
las seis coordenadas de la frontera regional.}
\label{act21:tab:terminal-bandas-testigo}
\end{table}

La tabla siguiente despliega todas las incidencias comparativas del
registro sobre los doce bloques de la segmentación considerada. Las
letras $r$ y $d$ designan respectivamente la carga visible y la carga
dual. La columna $j$ conserva el índice de posición del testigo histórico;
la operación de entrega al selector olvida este índice después de formar
el conjunto de palabras.

\begin{table}[htbp]
\centering\small
\begin{tabular}{rrcrrc c}
\toprule
$j$&$t$&carga&$\delta$&$\epsilon$&$(\eta_0\eta_1\eta_2)$&$L_{t,c,\delta,\epsilon}$\\
\midrule
4&11&$d$&0& 1&212&021101\\
5&67&$d$&1&$-1$&211&002001\\
5&67&$d$&2& 1&211&002001\\
5&67&$r$&1&$-1$&211&002001\\
5&67&$r$&2& 1&211&002001\\
6&81&$d$&0& 1&222&020111\\
6&81&$r$&2& 1&222&020111\\
7&34&$r$&1&$-1$&111&200110\\
7&75&$d$&1&$-1$&211&200110\\
7&75&$r$&2&$-1$&211&200110\\
8&90&$r$&0& 1&121&121012\\
8&90&$r$&1&$-1$&121&121012\\
9&49&$d$&0&$-1$&121&010122\\
11&11&$d$&2&$-1$&211&221110\\
11&37&$d$&0&$-1$&121&221110\\
12&8&$d$&0&$-1$&111&222110\\
12&8&$r$&1&$-1$&111&222110\\
12&58&$d$&1&$-1$&111&222110\\
12&58&$r$&0&$-1$&111&222110\\
\bottomrule
\end{tabular}
\caption{Incidencias fase--carga completas sobre la segmentación terminal
en las cien fronteras del registro. Las coincidencias de palabras entre
filas conservan sus tiempos, orientaciones y tipos de carga.}
\label{act21:tab:terminal-incidencias-completas}
\end{table}

El décimo bloque dispone de un testigo afín. En $t=30$ se tiene
$d_{30}=(0,1,2)$ y $\vartheta_{30}=2$, de donde
\[
 (d_{30,0}+\vartheta_{30},d_{30,1}+\vartheta_{30},
 d_{30,2}+\vartheta_{30})=(2,0,1)
\]
y, usando las bandas de la tabla~\ref{act21:tab:terminal-bandas-testigo},
\[
 L^{\mathrm{af}}_{30}
 =2(210112)+0(110202)+(211110)
 =(001001)\quad\text{en }\mathbb F_3^6.
\]
La igualdad es coordenada a coordenada: antes de reducir módulo tres,
el vector de la suma es $(6,3,1,3,3,4)$.

\begin{proposition}[Recuperación íntegra del repertorio por incidencias]
\label{act21:prop:repertorio-terminal-ocho-incidencias}
Considérese el registro de incidencias comparativas de la
tabla~\ref{act21:tab:terminal-incidencias-completas}, aumentado con el testigo
afín $(j,t,L)=(10,30,001001)$. Para cada posición $5\leq j\leq12$, el
conjunto de palabras de salida de sus incidencias es unitario. La
extracción de esas ocho palabras y el olvido posterior de su orden
producen exactamente
\[
 \mathcal S_{\mathrm{term}}=
 \{002001,020111,200110,121012,010122,001001,221110,222110\}.
\]
Su codificación como enteros de base tres es
\[
 \{28,55,98,175,437,498,687,714\}.
\]
\end{proposition}
\begin{proof}
La multiplicidad de incidencias por posición $j=5,6,7,8,9,10,11,12$
es, respectivamente,
\[
 (4,2,3,2,1,1,2,4).
\]
En cada posición, todas las filas imprimen la misma palabra. Cada fila
se verifica sustituyendo sus tres bandas y los tres coeficientes en la
fórmula de $L$; el décimo caso se ha calculado explícitamente mediante
$L^{\mathrm{af}}_{30}$. Por ejemplo, en $t=67$ y con coeficientes
$(2,1,1)$ se obtiene
\[
 2(010012)+(001220)+(011120)=(002001)\pmod3.
\]
Los ocho resultados son distintos como palabras de longitud seis. Su
codificación $\sum_{k=1}^{6}w_k3^{6-k}$ da, en el orden de las
posiciones, $(55,175,498,437,98,28,714,687)$. Al olvidar el orden se
obtiene el conjunto de enteros enunciado. Así, la extracción preserva
cada palabra completa, incluidas sus posiciones nulas, y entrega ocho
elementos distintos al selector.
\end{proof}

La prueba anterior es una recuperación exacta desde incidencias
identificadas. Su dato de entrada comprende el registro de incidencias
con sus posiciones y testigos. La selección orbital posterior recibe
únicamente $\mathcal S_{\mathrm{term}}$ y el descriptor regional inicial;
recorre las $8!$ ordenaciones y determina su compatibilidad con las
fronteras funcionales y las incidencias orientadas de la rueda
dodecafásica. Ambas operaciones quedan así contiguas, con dominios
explícitos: procedencia temporal de las ocho palabras y determinación
constructiva de su orden.

El cuadro comparativo contiene las diecinueve incidencias que su fórmula
produce sobre la segmentación de doce bloques en los cien tiempos
declarados. Esa exhaustividad finita se comprueba recorriendo
$100\cdot2\cdot3\cdot2=1200$ evaluaciones. El lector afín amplía la
cobertura de las posiciones terminales de siete a ocho. Este recuento
caracteriza el repertorio local de lectores que acaba de definirse;
las demás operaciones del TPK conservan los dominios y funciones
constructivas establecidos en sus respectivas derivaciones.

\subsection{Lectores comparativos, lector afín e incidencia temporal}

\subsubsection{Transformación de Witt del selector}

El selector utiliza la coordenatización de seis coordenadas
\[
 A_{\rm ter}=
 \begin{pmatrix}
 0&1&1&1&1&1\\
 1&0&1&2&2&1\\
 1&1&0&1&2&2\\
 1&2&1&0&1&2\\
 1&2&2&1&0&1\\
 1&1&2&2&1&0
 \end{pmatrix}.
\]
Se define $W(b)=bA_{\rm ter}$ en $\mathbb F_3^6$, así como las cargas
\[
 \eta(b)=\sum_{j=0}^5 b_j\pmod3,
 \qquad\eta_{\rm ter}^\vee(b)=\eta(W(b)).
\]
La coordenatización del descriptor y la del selector se relacionan mediante la
permutación $\sigma=(0,1,2,4,5,3)$, escrita como lista de imágenes:
\begin{equation}
 (A_{\rm ter})_{ij}=(A_{\rm reg})_{\sigma(i),\sigma(j)}.
 \label{act21:eq:kdes-cambio-carta}
\end{equation}

\begin{lemma}[Compatibilidad del funcional de carga]
Para toda banda $b$, se cumple
\[
 \eta_{\rm ter}^\vee(b)=\eta^\vee(b).
\]
\end{lemma}
\begin{proof}
La suma entera de cada fila correspondiente de ambas matrices coincide:
en la primera fila es cinco y en las otras cinco es siete. Por
distributividad,
\[
 \sum_j\sum_i b_i(A_{\rm ter})_{ij}
 =\sum_i b_i\sum_j(A_{\rm ter})_{ij}
 =\sum_i b_i\sum_j(A_{\rm reg})_{ij}.
\]
La reducción módulo tres da la igualdad de cargas. La identidad de
matrices \eqref{act21:eq:kdes-cambio-carta} se verifica por sustitución de sus
treinta y seis entradas y conserva, además, el cambio de coordenadas.
\end{proof}

La igualdad del funcional total de carga permite comparar sus
representaciones. Las transformaciones vectoriales siguen identificadas
por sus matrices y por la permutación explícita; la prueba anterior
conserva esa información de coordenadas.

\subsubsection{Doce lectores comparativos por instante}

Se define la fase temporal
\[
 \theta(t)=(t+2)\bmod3.
\]
Esta escritura realiza $(t-1)\bmod3$ también para $t=0$, donde la fase
vale dos. Para cada tipo de carga
$\nu\in\{\eta,\eta_{\rm ter}^\vee\}$, desplazamiento
$o\in\{0,1,2\}$ y orientación $s\in\{1,2\}$, se toman
\[
 \epsilon_\chi^{\nu,o,s}(t)=s
 \begin{cases}
 1,&\nu(B_\chi(t))=\theta(t)+o\pmod3,\\
 2,&\nu(B_\chi(t))\ne\theta(t)+o\pmod3.
 \end{cases}
\]
La palabra comparativa correspondiente es
\begin{equation}
 C_{\nu,o,s}(t)=
 \sum_{\chi\in\{\pi,e,\varphi\}}
 \epsilon_\chi^{\nu,o,s}(t)B_\chi(t)
 \quad\text{en }\mathbb F_3^6.
 \label{act21:eq:kdes-comparativo}
\end{equation}
Los parámetros generan doce representaciones etiquetadas por fila. Dos
etiquetas pueden producir la misma palabra; los testigos temporales y
las etiquetas permanecen disponibles en el calendario completo.

\subsubsection{Lector afín de fase y carga}

El lector afín emplea coeficientes
\[
 a_\chi^{\rm af}(t)=\eta_{\rm ter}^\vee(B_\chi(t))+\theta(t)
 \pmod3,
\]
y produce
\begin{equation}
 F(t)=\sum_{\chi\in\{\pi,e,\varphi\}}
 a_\chi^{\rm af}(t)B_\chi(t)
 \quad\text{en }\mathbb F_3^6.
 \label{act21:eq:kdes-afin}
\end{equation}
La regla comparativa distingue coincidencia y diferencia de cargas;
la afín suma carga y fase antes de combinar las bandas. Esta diferencia
operativa es la que utiliza la selección heterotípica.

\subsubsection{Incidencia palabra--residuo}

Sea $R$ la rotación cíclica de las seis posiciones de una banda. En la
fila de tiempo $t$ se forma el repertorio residual
\[
 \mathcal R(t)=\{R^jv:
      v\in\{q(t),a(t),c(t),\bar w(t)\},\ 0\leq j<6\}.
\]
Las relaciones de incidencia temporal son
\begin{align}
 \mathcal A_{\rm cmp}
 &=\{(C_{\nu,o,s}(t),u):
      0\leq t<100,\ u\in\mathcal R(t),\ \nu,o,s\text{ admisibles}\},\\
 \mathcal A_{\rm af}
 &=\{(F(t),u):0\leq t<100,\ u\in\mathcal R(t)\}.
 \label{act21:eq:kdes-atlas}
\end{align}
Cada par relaciona una palabra leída con un residuo de la misma fila.
El tiempo se cuantifica de manera común en ambas coordenadas del par.
En el calendario empleado, los tiempos $0,\ldots,99$ son distintos;
identificar una misma fila e identificar un mismo tiempo son, por ello,
operaciones equivalentes.

\begin{proposition}[Compresión existencial del repertorio temporal]
La función
\[
 M(v,u)=\mathbf1_{(v,u)\in\mathcal A_{\rm cmp}}
        +2\mathbf1_{(v,u)\in\mathcal A_{\rm af}}
        \in\{0,1,2,3\}
\]
conserva exactamente la existencia de cada tipo de lectura para el par
$(v,u)$.
\end{proposition}
\begin{proof}
El primer bit vale uno exactamente cuando existe un tiempo, un campo
residual, una rotación y una etiqueta comparativa que producen el par.
El segundo bit expresa la condición correspondiente al lector afín.
La unión por disyunción de bits conserva cada condición de existencia
al recorrer todas las filas. Un par puede admitir ambos tipos, y en ese
caso la máscara vale tres.
\end{proof}

La función $M$ es un lector de existencia. Los tiempos testigos, sus
multiplicidades y las etiquetas individuales permanecen en
\eqref{act21:eq:kdes-atlas} y en el calendario que las produce. La
implementación finita emplea $M$ para decidir heterotipia, mientras que
el archivo genealógico conserva los datos de procedencia más ricos.

\subsection{Enumeración de cilindros y candidatos decimales}

\subsubsection{Concatenación de setenta y dos y setenta y ocho trits}

Para una ordenación $\tau=(s_1,\ldots,s_8)$ de
$\mathcal S_{\rm term}$ se define
\[
 P_{72}(\tau)=[\Omega_{24}\Vert s_1\Vert\cdots\Vert s_8]_3.
\]
En forma recursiva, se parte de $p_0=[\Omega_{24}]_3$ y se calcula
$p_i=729p_{i-1}+[s_i]_3$ para $1\leq i\leq8$. Esta recurrencia
conserva los ceros iniciales de cada bloque por la multiplicación
explícita por $729$. Para cada posición $b$ del panel funcional,
\[
 P_{78}(\tau,b)=729P_{72}(\tau)+b.
\]
Las longitudes son, respectivamente, $24+8\cdot6=72$ y $72+6=78$.
El cilindro ternario correspondiente es
\begin{equation}
 I_{\tau,b}=
 \left[\frac{P_{78}(\tau,b)}{3^{78}},
       \frac{P_{78}(\tau,b)+1}{3^{78}}\right).
 \label{act21:eq:kdes-cilindro}
\end{equation}

\subsubsection{Intersección exacta con la rejilla de doce tríadas}

La lectura decimal utiliza doce tríadas, es decir, la rejilla
$N/10^{36}$. Se abrevia $D=10^{36}$ y $Q=3^{78}$.

\begin{lemma}[Extremos enteros del cilindro]
Para un prefijo entero $p$ de longitud setenta y ocho, los numeradores
de los puntos de la rejilla que pertenecen a su cilindro son
\begin{equation}
 \left\{N\in\mathbb Z:
 \left\lceil\frac{pD}{Q}\right\rceil
 \leq N<
 \left\lceil\frac{(p+1)D}{Q}\right\rceil\right\}.
 \label{act21:eq:kdes-rejilla}
\end{equation}
Cada cilindro contiene a lo sumo un punto de esta rejilla.
\end{lemma}
\begin{proof}
La condición $N/D\in[p/Q,(p+1)/Q)$ equivale a
$pD/Q\leq N<(p+1)D/Q$. Para $N$ entero, la primera desigualdad
equivale a $\lceil pD/Q\rceil\leq N$; la segunda equivale a
$N<\lceil(p+1)D/Q\rceil$. La longitud del intervalo en coordenada
$N$ es $D/Q<1$, pues $10^{36}<3^{78}$. Por ello puede contener,
como máximo, un entero.
\end{proof}

Los techos se evalúan por división entera:
$\lceil a/Q\rceil=\lfloor(a+Q-1)/Q\rfloor$ para $a\geq0$.
Toda la enumeración utiliza enteros exactos; la elección de puntos de
frontera queda fijada por el extremo derecho semiabierto.

Se conserva primero la lista indexada de incidencias
\[
 \mathcal C_{\rm ind}=
 \{(\tau,j,N):N/D\in I_{\tau,\mathcal P_9(j)}\},
 \qquad1\leq j\leq9,
\]
y se obtiene después el conjunto de numeradores
\[
 \mathcal C=\{N:\exists\tau,j\ (\tau,j,N)\in\mathcal C_{\rm ind}\}.
\]
El paso a $\mathcal C$ elimina repeticiones de un numerador, conservadas
en $\mathcal C_{\rm ind}$ con su procedencia.

\begin{lemma}[Recuperación de la ordenación desde el punto del cilindro]
Si $(\tau,j,N)\in\mathcal C_{\rm ind}$, el bloque número $i$ de la
palabra de setenta y ocho trits se recupera mediante
\[
 \left\lfloor\frac{N\,729^i}{D}\right\rfloor\bmod729,
 \qquad1\leq i\leq13.
\]
En particular, los bloques quinto a duodécimo recuperan $\tau$.
\end{lemma}
\begin{proof}
La pertenencia al cilindro significa
$p\leq QN/D<p+1$, donde $p=P_{78}(\tau,\mathcal P_9(j))$.
Por tanto, $\lfloor QN/D\rfloor=p$. La extracción de los bloques
anteriores coincide con la división sucesiva del prefijo $p$ por
potencias de $729$: truncar una expansión posicional antes de extraer
un bloque conserva dicho bloque. La fórmula escrita realiza
exactamente esa extracción. Los cuatro primeros bloques forman
$\Omega_{24}$ y los ocho siguientes son la ordenación $\tau$.
\end{proof}

\begin{proposition}[Censo exacto de candidatos]
Para $\Omega_{24}$, $\mathcal S_{\rm term}$ y $\mathcal P_9$
definidos anteriormente,
\[
 \#\operatorname{Ord}(\mathcal S_{\rm term})=40320,
 \qquad\#\mathcal C_{\rm ind}=21902,
 \qquad\#\mathcal C=19446.
\]
\end{proposition}
\begin{proof}
Los ocho bloques son distintos, por lo que sus ordenaciones son las
$8!=40320$ permutaciones. Para cada permutación y cada una de las
nueve posiciones del panel se calcula $p=P_{78}(\tau,b)$, se aplican
los extremos \eqref{act21:eq:kdes-rejilla} y se añade su único entero cuando
el intervalo entero es no vacío. La suma de los cardinales de esos
intervalos es $21902$. La unión de sus enteros contiene $19446$
elementos distintos. La evaluación exhaustiva de estas operaciones
está certificada por los teoremas
\texttt{indexed\_cylinder\_count} y \texttt{candidate\_count} del
selector finito. El algoritmo realiza exactamente la enumeración
descrita: permutaciones, nueve incidencias del panel, límites por
división entera y eliminación de duplicados.
\end{proof}

Los números $21902$ y $19446$ cuentan objetos diferentes. El primero
cuenta incidencias indexadas que contienen un punto decimal; el segundo
cuenta numeradores distintos. El número total de consultas de cilindro
es $40320\cdot9=362880$, incluidos los cilindros cuya intersección con
la rejilla es vacía.

\subsection{Incidencia excepcional y condición de heterotipia}

\subsubsection{Cara excepcional producida por el panel}

Se eleva una banda a doce coordenadas mediante
\[
 \Gamma(b)=(b,W(b))\in\mathbb F_3^{12}.
\]
El soporte de una palabra es el conjunto de posiciones de sus
coordenadas no nulas, con índices humanos $1,\ldots,12$. Las primeras
bandas regionales del panel producen
\begin{equation}
 \mathcal B=\operatorname{supp}\Gamma(B_\pi(0))
          \cap\operatorname{supp}\Gamma(B_e(0))
          =\{4,6,9,10\}.
 \label{act21:eq:kdes-cara}
\end{equation}
La operación utiliza las palabras regionales $103$ y $523$ y la
matriz de Witt del selector. La cara se calcula antes de consultar las
coordenadas de los candidatos.

Para $N\in\mathcal C$, sus doce tríadas son
\begin{equation}
 K_i(N)=\left\lfloor\frac{N}{1000^{12-i}}\right\rfloor\bmod1000,
 \qquad1\leq i\leq12.
 \label{act21:eq:kdes-triadas}
\end{equation}
El soporte de corona es
\[
 \mathcal H(N)=\{i:K_i(N)\geq729\},
 \qquad
 \mathcal C_{\mathcal B}=\{N\in\mathcal C:\mathcal H(N)=\mathcal B\}.
\]
La frontera $729=3^6$ relaciona la banda de seis trits con la tríada
decimal. La evaluación exhaustiva de la igualdad de soportes da
\begin{equation}
 \#\mathcal C_{\mathcal B}=79.
 \label{act21:eq:kdes-79}
\end{equation}
El soporte conserva posiciones, mientras que el registro completo
conserva también las doce magnitudes. La etapa siguiente utiliza esa
información posicional junto con relaciones entre magnitudes.

\subsubsection{Palabras ternarias y diferencias orientadas del candidato}

La $i$-ésima banda ternaria del punto $N/D$ se obtiene por
\begin{equation}
 T_i(N)=\left\lfloor\frac{N\,729^i}{D}\right\rfloor\bmod729,
 \qquad1\leq i\leq13.
 \label{act21:eq:kdes-tercandidato}
\end{equation}
La arista orientada $a\to b$ tiene residuo
\begin{equation}
 R_{a,b}(N)=(K_a(N)-K_b(N))\bmod729.
 \label{act21:eq:kdes-resarista}
\end{equation}
La resta se realiza en $\mathbb Z$ antes de la reducción. Para esta
arista se consulta en las relaciones \eqref{act21:eq:kdes-atlas} el par
\[
 \bigl(\operatorname{enc}_6^{-1}(T_b(N)),
  \operatorname{enc}_6^{-1}(R_{a,b}(N))\bigr).
\]
Es, por tanto, una consulta
conjunta a la palabra de destino y a la diferencia orientada de las
tríadas.

Se escribe $\mathrm{Cmp}(a,b;N)$ cuando el par pertenece a
$\mathcal A_{\rm cmp}$ y $\mathrm{Af}(a,b;N)$ cuando pertenece a
$\mathcal A_{\rm af}$.

\subsubsection{Determinación de la órbita transversal}

El descriptor de veinticuatro trits determina las tres primeras
tríadas decimales antes de las filtraciones terminales. En efecto, su
cilindro satisface la inclusión exacta
\[
 \left[\frac{66241910521}{3^{24}},
       \frac{66241910522}{3^{24}}\right)
 \subseteq
 \left[\frac{234543140}{10^9},
       \frac{234543141}{10^9}\right).
\]
Las dos desigualdades de extremos se verifican por productos cruzados
enteros. Por tanto, todo candidato comparte las tres posiciones de
referencia $1,2,3$ y sus tríadas decimales $(234,543,140)$.

La traslación de paso cuatro sobre $\mathbb Z/12\mathbb Z$ posee las
cuatro órbitas
\[
 (1,5,9),\qquad(2,6,10),\qquad(3,7,11),\qquad(4,8,12).
\]
Las tres primeras contienen, respectivamente, una de las posiciones de
referencia; la cuarta es la única disjunta de $\{1,2,3\}$. El bosque
orientado de los pasos tres y cuatro se fija mediante las aristas
\[
 \begin{array}{lll}
 1\to4,&1\to5,&2\to6,\\
 3\to7,&4\to8,&5\to9,\\
 6\to10,&7\to11,&8\to12.
 \end{array}
\]
Su primera arista tiene paso tres y las restantes paso cuatro. Las dos
aristas internas de la órbita distinguida son precisamente
$4\to8$ y $8\to12$. Así queda fijado el lugar de la condición de
lectura por la geometría de las posiciones y por el prefijo inicial,
antes de examinar los candidatos de la cara excepcional.

\subsubsection{Condición terminal de heterotipia}

La condición terminal es
\begin{equation}
 \begin{split}
 \mathcal T(N)={}&
 [\mathrm{Cmp}(4,8;N)\wedge\mathrm{Af}(8,12;N)]\\
 &\vee[\mathrm{Af}(4,8;N)\wedge\mathrm{Cmp}(8,12;N)].
 \end{split}
 \label{act21:eq:kdes-heterotipia}
\end{equation}
Se admiten las dos asignaciones de tipos a las dos aristas. La
selección exige que exista una realización con tipos distintos; las
etiquetas individuales y sus tiempos quedan disponibles en los
testigos de incidencia.

\begin{theorem}[Selección terminal del registro de acoplamiento]
Sobre el dominio de candidatos producido por
$\Omega_{24}$, $\mathcal S_{\rm term}$, $\mathcal P_9$ y el calendario
regional de cien filas, existe un único numerador que satisface
simultáneamente la incidencia excepcional y la heterotipia:
\[
 \{N\in\mathcal C_{\mathcal B}:\mathcal T(N)\}
 =\{234543140729659824621058914794146601\}.
\]
Su registro de doce tríadas es
\begin{equation}
 K=(234,543,140,729,659,824,621,058,914,794,146,601).
 \label{act21:eq:kdes-kfinal}
\end{equation}
\end{theorem}
\begin{proof}
El censo precedente produce los $79$ elementos de
$\mathcal C_{\mathcal B}$. Para cada uno se extraen sus doce
coordenadas por \eqref{act21:eq:kdes-triadas}; se calculan las bandas de
destino $T_8,T_{12}$ y los residuos $R_{4,8},R_{8,12}$; se consultan
las dos máscaras del repertorio temporal y se aplica la disyunción
\eqref{act21:eq:kdes-heterotipia}. La lista resultante contiene exactamente
el entero indicado, según la evaluación exhaustiva certificada por
\texttt{terminal\_selection}. La igualdad de lista con un singleton
implica existencia y unicidad en el dominio declarado. La extracción
base mil de ese entero da \eqref{act21:eq:kdes-kfinal}.

El procedimiento de selección recibe los productores y las reglas
enumeradas. El entero final interviene en el enunciado de la igualdad
comprobada, después de ejecutar la selección. Esta posición lógica
permite cotejar el resultado sin utilizarlo como criterio de elección.
\end{proof}

\subsubsection{Testigo aritmético de pertenencia a un cilindro}

Una de las incidencias del numerador seleccionado utiliza la ordenación
\[
 (002001,020111,200110,121012,010122,001001,221110,222110)
\]
y la frontera $438=121020_3$. Para ella,
\[
 P_{72}=5283881584882673136514102926090957,
\]
\[
 P_{78}=3851949675379468716518781033120308091.
\]
La fórmula de extremos enteros produce
\begin{align*}
 \left\lceil\frac{P_{78}10^{36}}{3^{78}}\right\rceil
 &=234543140729659824621058914794146601,\\
 \left\lceil\frac{(P_{78}+1)10^{36}}{3^{78}}\right\rceil
 &=234543140729659824621058914794146602.
\end{align*}
El único entero del intervalo es, por tanto, $N_K$. Las trece bandas
ternarias de su punto decimal son
\[
 \begin{array}{rrrr}
 020022&222111&211201&021101\\
 002001&020111&200110&121012\\
 010122&001001&221110&222110\\
 121020&&&
 \end{array}
\]
Este cálculo exhibe un testigo concreto de pertenencia al dominio
enumerado. La unicidad del registro procede del examen exhaustivo de
todos los candidatos, además de este testigo particular.

\subsection{Composición regional y continuidad de las realizaciones}

\subsubsection{Composición de los generadores regionales}

Sea $\operatorname{Sel}(\Omega,\mathcal S,\mathcal P,\mathcal N)$ el
selector definido por la enumeración y los dos filtros anteriores.
Los productores regionales satisfacen las igualdades
\[
 \Omega_{\rm reg}=\Omega_{24},\qquad
 \mathcal P_{\rm reg}=\mathcal P_9,\qquad
 \mathcal N_{\rm reg}=\mathcal N_{100},
\]
donde $\mathcal N_{\rm reg}$ se calcula por extracción de bandas y
lectura de firmas. Por congruencia de funciones,
\begin{equation}
 \begin{split}
 &\operatorname{Sel}(\Omega_{\rm reg},\mathcal S_{\rm term},
                    \mathcal P_{\rm reg},\mathcal N_{\rm reg})\\
 &\hspace{10mm}=
 \operatorname{Sel}(\Omega_{24},\mathcal S_{\rm term},
                    \mathcal P_9,\mathcal N_{100})=\{N_K\}.
 \end{split}
 \label{act21:eq:kdes-composicion}
\end{equation}
La igualdad reemplaza tres interfaces documentales por su construcción
regional efectiva. El repertorio terminal mantiene la procedencia
especificada en \eqref{act21:eq:kdes-sterm}.

El registro se forma entonces por las doce divisiones euclídeas de
\eqref{act21:eq:kdes-triadas}. Cada coordenada pertenece a
$\{0,\ldots,999\}$ y la longitud es doce por construcción. La
implementación que toma el primer elemento de la lista seleccionada
está justificada por la igualdad con el singleton: el valor por defecto
de una operación informática sobre listas vacías nunca interviene en
este resultado.

\subsubsection{Relación con las incidencias firmadas y la inversión}

La construcción presente proporciona una selección prospectiva en un
dominio finito explícito. El desarrollo ulterior del registro
dodecafásico produce sus canales orientados a partir del estado
terminal enriquecido y demuestra una inversión integral mediante
Hadamard. Aquella inversión recupera un registro desde canales
compatibles; aquí se selecciona un registro a partir de sus
antecedentes regionales e incidenciales. La igualdad de sus resultados
permite componer las dos realizaciones manteniendo el origen de cada
operación.

La unicidad de selección y la unicidad de inversión son propiedades
distintas. La primera cuantifica los candidatos construidos en este
capítulo. La segunda cuantifica los registros que comparten una firma
compatible. Su conjunción justifica que el registro seleccionado pueda
transportarse a la realización firmada y recuperarse desde ella, sin
convertir una recuperación retrospectiva en el productor del registro.

\subsubsection{Relación entre los censos terminales}

La secuencia
\[
 40320\ \longrightarrow\ 21902\ \longrightarrow\ 19446
 \ \longrightarrow\ 79\ \longrightarrow\ 1
\]
registra, sucesivamente, ordenaciones, incidencias indexadas con la
rejilla, numeradores distintos, candidatos de la cara excepcional y
supervivientes heterotípicos. La notación de flechas indica el orden de
las operaciones; la segunda cifra cuenta una clase de objetos
distinta de la primera.

El censo de catálogos regionales
$196\cdot197\cdot196\longrightarrow331\longrightarrow2
\longrightarrow1$, desarrollado en su lugar correspondiente, parte de
otra población: bloques regionales de tres catálogos y condiciones
terminales de Gram, ocupación, distancia y orientación. Ambos censos
conservan sus dominios. Su convergencia en una misma representación se
expresa mediante la identificación de sus lectores y de sus salidas;
las cifras intermedias describen operaciones diferentes y permanecen
documentadas por separado.

\subsubsection{Composición regional de la constante de estructura fina}

La salida $K$ interviene a continuación en la composición regional
ordenada y en el acarreo de base mil. Su empleo posterior conserva las
doce posiciones y la frontera de lectura. El acoplamiento con la raíz
analítica y las representaciones a profundidad arbitraria utiliza esa
misma salida, junto con las coordenadas regionales ya generadas y las
reglas de la coordenatización analítica. Esta dependencia prepara el capítulo
dedicado a las dos realizaciones correlacionadas de la constante de
estructura fina.

El alcance de este capítulo es preciso: producción del descriptor y
de sus antecedentes regionales, determinación del dominio finito,
selección única sobre dicho dominio y composición material con los
productores. La recurrencia analítica de todos los grados y la
compatibilidad de los acarreos infinitos se desarrollan después con sus
demostraciones propias.

\subsection{Certificación ejecutable y trazabilidad de las operaciones}

El desarrollo formal conserva la correspondencia entre las
definiciones matemáticas y sus realizaciones ejecutables. Las pruebas
de firmas y de reloj se encuentran en
\texttt{N69RegionalSignature.lean}; la extracción de bandas y la
reproducción de las cien filas, en
\texttt{GeneratedN69Rows.lean}; el criterio mínimo y la construcción
de $\Omega_{24}$, en \texttt{RegionalW24.lean}. La compatibilidad de
las dos matrices de Witt está reunida en
\texttt{WittReaderCompatibility.lean}.

Los lectores temporales, las rotaciones, los residuos orientados y la
condición heterotípica se formalizan en
\texttt{TerminalReaders.lean}. La enumeración y los censos se
formalizan en \texttt{Terminal}\allowbreak\texttt{Orbital}\allowbreak\texttt{Selection.lean}. El panel
regional se enlaza mediante \texttt{Regional}\allowbreak\texttt{Panel}\allowbreak\texttt{Bridge.lean}, y la
composición \eqref{act21:eq:kdes-composicion} queda expresada en
\texttt{Selected}\allowbreak\texttt{From}\allowbreak%
\texttt{Regional}\allowbreak\texttt{Inputs.lean}.

Las comprobaciones finitas de enumeración emplean la evaluación nativa
de Lean; sus informes registran \texttt{Lean.ofReduceBool}, además
de los axiomas lógicos habituales de la biblioteca. La reconstrucción
regional del descriptor utiliza pruebas del núcleo y reducción
ordinaria. Esta distinción describe los mecanismos de comprobación
realmente empleados. Los recibos de compilación acreditan sus módulos
y dependencias declarados; las pruebas matemáticas conservan los
dominios, productores e hipótesis que se han explicitado en el cuerpo
del capítulo.
